% main_clean.tex -- the manuscript, and the only hand-edited paper file. BLUEPRINT.md is its design.
% Order of the contributions, everywhere: (1) a common measurement model, (2) a taxonomy by three questions,
% (3) a review of each class with one representative method, (4) every representative method on the same scans.
% Math (2026-11-10): as little as the argument needs, in the notation of the papers each class comes from, with
% the source of the notation stated. No new symbols. The simplest form of an algorithm is the one shown.
% Conventions: every number is \V{key} from values.tex (make_values.py, from the notebook's results/*.csv);
% every figure is written by the notebook; the reference list is built by make_references.py from
% LITERATURE.md; \outline{} is the plan of a section not yet written and \pending{} an open item, both red.
\documentclass[10pt,letterpaper]{article}

\usepackage[utf8]{inputenc}
\usepackage[T1]{fontenc}
\usepackage{newtxtext}
\usepackage[letterpaper,textwidth=6.5in,textheight=9in,top=1in,footskip=30pt]{geometry}
\usepackage{amsmath,amssymb,bm,mathtools}
\usepackage{microtype}
\usepackage[table]{xcolor}
\usepackage{graphicx}
\graphicspath{{figures/}}
\usepackage{booktabs}
\usepackage{siunitx}
\usepackage{array}
\usepackage{caption}
\captionsetup{labelfont=bf}
\usepackage[hidelinks]{hyperref}
% In-text references to floats, sections, and equations: the whole "Fig. 2" is bold and is the link.
\newcommand{\xref}[3][]{\hyperref[#3]{\textbf{#2~\ref*{#3}#1}}}
\newcommand{\xeq}[1]{\hyperref[#1]{\textbf{Eq.~(\ref*{#1})}}}
\usepackage{cite}

% Paragraphs: no indent, a small gap between them.
\setlength{\parindent}{0pt}
\setlength{\parskip}{5.5pt}

% Section headings: bold, modest sizes.
\makeatletter
\renewcommand\section{\@startsection{section}{1}{\z@}%
  {-2.0ex \@plus -0.5ex \@minus -0.2ex}{1.5ex \@plus 0.3ex}{\large\bfseries}}
\renewcommand\subsection{\@startsection{subsection}{2}{\z@}%
  {-1.8ex \@plus -0.5ex \@minus -0.2ex}{0.8ex \@plus 0.2ex}{\normalsize\bfseries}}

% Title block: small-caps title, no rules, centered authors.
\renewcommand{\@maketitle}{%
  \begin{center}
    {\LARGE\scshape \@title\par}
    \vskip 0.3in
    \@author
  \end{center}
  \vskip 0.15in}
\makeatother

% Floats: let figures and tables sit close to their first mention.
\renewcommand{\topfraction}{0.9}
\renewcommand{\bottomfraction}{0.8}
\renewcommand{\textfraction}{0.07}
\renewcommand{\floatpagefraction}{0.8}

% ---------- Placeholders ----------
\newcommand{\pending}[1]{{\color{red!75!black}#1}}
\newcommand{\outline}[1]{{\color{red!75!black}\small\textit{To write.} #1}}
% a figure the notebook has not written yet prints as a red box instead of stopping the build
\newcommand{\figureorpending}[2]{\IfFileExists{figures/#1}{\includegraphics[width=\textwidth,height=0.8\textheight,keepaspectratio]{#1}}{\fbox{\parbox[c][4cm][c]{0.9\textwidth}{\centering\color{red!75!black}\small [\detokenize{#1}: #2]}}}}
\makeatletter
\newcommand{\V}[1]{\@ifundefined{v@#1}{{\color{red!75!black}[\detokenize{#1}]}}{\csname v@#1\endcsname}}
\makeatother
\IfFileExists{values.tex}{% Generated by paper/make_values.py from results/*.csv. Do not edit.
\expandafter\def\csname v@census-first-year\endcsname{1976}
\expandafter\def\csname v@census-last-all\endcsname{115}
\expandafter\def\csname v@census-last-learned\endcsname{17}
\expandafter\def\csname v@census-last-review\endcsname{16}
\expandafter\def\csname v@census-last-title\endcsname{19}
\expandafter\def\csname v@census-last-year\endcsname{2025}
\expandafter\def\csname v@census-run\endcsname{2026-11-10}
\expandafter\def\csname v@census-total-all\endcsname{921}
\expandafter\def\csname v@census-total-learned\endcsname{111}
\expandafter\def\csname v@census-total-review\endcsname{85}
\expandafter\def\csname v@census-total-title\endcsname{225}
\expandafter\def\csname v@mm-bound-iodine-best\endcsname{3.08}
\expandafter\def\csname v@mm-bound-iodine-no-tin\endcsname{3.77}
\expandafter\def\csname v@mm-correlation-best\endcsname{-0.932}
\expandafter\def\csname v@mm-determined-de-2\endcsname{2}
\expandafter\def\csname v@mm-determined-de-3\endcsname{2}
\expandafter\def\csname v@mm-determined-de-5\endcsname{2}
\expandafter\def\csname v@mm-determined-pc-2\endcsname{2}
\expandafter\def\csname v@mm-determined-pc-3\endcsname{3}
\expandafter\def\csname v@mm-determined-pc-5\endcsname{3}
\expandafter\def\csname v@mm-determined-se-2\endcsname{1}
\expandafter\def\csname v@mm-determined-se-3\endcsname{1}
\expandafter\def\csname v@mm-determined-se-5\endcsname{1}
\expandafter\def\csname v@mm-n-tissues\endcsname{10}
\expandafter\def\csname v@mm-residual-edges-2\endcsname{18}
\expandafter\def\csname v@mm-residual-edges-3\endcsname{7.8}
\expandafter\def\csname v@mm-residual-edges-4\endcsname{0.12}
\expandafter\def\csname v@mm-residual-tissues-1\endcsname{18}
\expandafter\def\csname v@mm-residual-tissues-2\endcsname{0.16}
\expandafter\def\csname v@mm-tin-best\endcsname{0.3}
\expandafter\def\csname v@ray-bias-image-max\endcsname{2.77}
\expandafter\def\csname v@ray-bias-image-max-percent\endcsname{14}
\expandafter\def\csname v@ray-bias-projection-max\endcsname{2.0e-13}
\expandafter\def\csname v@ray-iodine\endcsname{20}
\expandafter\def\csname v@ray-noise-draws\endcsname{2000}
\expandafter\def\csname v@ray-noise-image\endcsname{0.99}
\expandafter\def\csname v@ray-noise-projection\endcsname{1.02}
}{}

\title{Material Decomposition in CT:\\One Measurement Model from Two Basis Materials\\to Multi-Material and Learned Methods}
% Working title. Author list and affiliations to be set.
\author{\pending{Authors and affiliations to be set}}
\date{}

\begin{document}

\maketitle

\begin{center}
  {\large\scshape Abstract}
\end{center}
\begin{quote}
\outline{Written last. Problem (many classes of methods, reviewed by application or by scanner), gap (no
common framework and no common test), what was done (one measurement model, three questions, one
representative method for each class, every method on the same simulated scans in one open notebook), what was
found (numbers from the benchmark), why it matters.}

\medskip
\noindent\textbf{\textit{Keywords}}\quad material decomposition $\cdot$ dual-energy CT $\cdot$ photon-counting CT
$\cdot$ multi-material decomposition $\cdot$ single-energy CT $\cdot$ deep learning
\end{quote}

% =============================================================================
\section{Introduction}
\label{sec:introduction}
% =============================================================================

\outline{Paragraph 1: what material decomposition is and what it is used for (iodine maps, virtual
non-contrast and virtual monoenergetic images, tissue composition). Paragraph 2: the growth of the field. A
PubMed search of titles and abstracts found \V{census-total-all} papers on material decomposition in CT from
\V{census-first-year} to \V{census-last-year}, \V{census-last-all} of them in \V{census-last-year}, of which
\V{census-total-learned} name a learning method and \V{census-total-review} are reviews (searched on
\V{census-run}). Paragraph 3: the existing reviews, grouped by what they organize (application, scanner,
learning), each read first. Paragraph 4: the gap. No review states every class in one model, and no two
classes are tested on the same data. Paragraph 5: this work, with the four contributions in order.
Paragraph 6: scope (CT with x-ray tubes, clinical energies, K-edge imaging as one rung, no phase contrast).}

\begin{figure*}[!htbp]
\centering
\figureorpending{fig_census.pdf}{publications per year, to be written by the notebook from results/literature\_census.csv}
\caption{\textbf{Growth of the literature on material decomposition in CT.} \outline{PubMed records per year
that name material or basis decomposition with CT in the title or abstract, those that also name a learning
method, and those listed as reviews.}}
\label{fig:census}
\end{figure*}

% =============================================================================
\section{A common measurement model}
\label{sec:model}
% =============================================================================

Every method reviewed here starts from the same description of the measurement, so it is stated once. The
notation is that of the projection-domain literature~\cite{alvarez1976,lehmann1981,blackprojection}, extended
from two basis materials to $M$, and no new symbols are introduced. Over the range of diagnostic x-ray energies, the attenuation of a material can be well approximated as
the sum of a photoelectric and a Compton component, each with its own energy dependence~\cite{alvarez1976}.
Two materials with different energy dependence span the same two functions, so the attenuation can equally be
written as a sum over basis materials~\cite{lehmann1981},
\begin{equation}
\mu(\bm r,E)=\sum_{m=1}^{M}\rho_m(\bm r)\,f_m(E),
\label{eq:basis}
\end{equation}
where $f_m(E)$ is the mass attenuation coefficient of basis material $m$ at the energy $E$ and $\rho_m(\bm r)$
is its density at the position $\bm r$. For \V{mm-n-tissues} tissues and basis materials without a K-edge
between 20 and \SI{140}{keV}, two energy functions left \V{mm-residual-tissues-2}\% of the attenuation curves
unexplained. A material with a K-edge in the detected spectrum does not follow the two functions and can be
given a function of its own~\cite{roessl2007}. With iodine and gadolinium added, two functions left
\V{mm-residual-edges-2}\% unexplained, three \V{mm-residual-edges-3}\%, and four \V{mm-residual-edges-4}\%.

A CT scanner measures this attenuation with polychromatic beams. For one ray, measurement $k$ of $K$ has the
expected value
\begin{equation}
\lambda_k=\int S_k(E)\,\exp\!\Big[-\sum_{m=1}^{M}f_m(E)\,A_m\Big]\,dE,
\qquad A_m=\int_{\ell}\rho_m(\bm r)\,d\ell,
\label{eq:forward}
\end{equation}
where $A_m$ is the line integral of basis material $m$ along the ray $\ell$ and $S_k(E)$ is the spectral
response of the measurement, the expected number of photons detected per unit energy without the
object~\cite{blackprojection}. The spectral response includes the source spectrum, the filtration, and the
detector response, so the acquisition mechanism enters the model only through $S_k(E)$. A single-energy scan
has $K=1$. A dual-energy scan has $K=2$, whether its two responses come from two tube potentials, two sources,
or two detector layers~\cite{mccollough2015}. A photon-counting scan has one response for each energy
bin~\cite{rajendran2022,mccollough2023}.

Material decomposition is the inversion of \xeq{eq:forward} for the densities $\rho_m$. The methods differ in
three respects, which organize this review (\xref{Fig.}{fig:framework}).

\begin{figure*}[!htbp]
\centering
\figureorpending{fig_framework.pdf}{written by the notebook}
\caption{\textbf{One measurement model and three questions.} The materials of the object attenuate a
polychromatic beam, the scanner records $K$ spectral measurements of each ray, and a decomposition turns them
into material maps. The methods differ in where they invert the model, in what closes the system when a voxel
can hold more materials than the measurements determine, and in how they report the answer.}
\label{fig:framework}
\end{figure*}
% SOURCE: notebooks/material_decomposition.jl, "The Chain at a Glance" -> figures/fig_framework.pdf

The first is where the model is inverted. Projection-domain methods invert \xeq{eq:forward} for each ray and
then reconstruct the line integrals $A_m$ into basis images~\cite{alvarez1976}. They keep the polychromatic
model, and they need the $K$ measurements along the same ray path. Image-domain methods reconstruct the $K$
measurements first and then solve a linear system in each voxel, which is simple to implement and applies to
any acquisition but can retain beam-hardening artifacts unless they are corrected~\cite{maass2009}. One-step
methods estimate the basis images directly from all measurements, which keeps the polychromatic model and can
add spatial regularization at a higher computational cost~\cite{long2014}.

The second is what closes the system when a voxel can hold more materials than its measurements determine.
$K$ measurements can determine at most $K$ amounts in a voxel, and no more than the number of energy functions
that its materials span. With a third condition, such as the
conservation of volume or mass, two measurements can be used to estimate three materials~\cite{liu2009}, and
one measurement two~\cite{xue2021}. Beyond this count, the information has to come from outside the voxel's own
measurements. Multi-material decomposition restricts each voxel to one triplet from a predefined
list~\cite{mendonca2014}, statistical methods add a penalty on the material images~\cite{long2014,xue2017},
and learned methods take the missing information from training data~\cite{bousse2024review}.
\xref{Table}{tab:ladder} orders the classes by this count. For a ray through \SI{20}{cm} of water and ideal
counting detectors, a dual-energy acquisition determined \V{mm-determined-de-2} of two, \V{mm-determined-de-3}
of three, and \V{mm-determined-de-5} of five material amounts, and four energy bins of one beam determined
\V{mm-determined-pc-2}, \V{mm-determined-pc-3}, and \V{mm-determined-pc-5}, where an amount counts as
determined when its singular value is above 1\% of the largest.

\begin{table*}[!htbp]
\centering
\caption{\textbf{The classes of material decomposition, ordered by what closes the system.} $K$ is the number
of spectral measurements of a ray or voxel.}
\label{tab:ladder}
\small
\begin{tabular}{@{}llllll@{}}
\toprule
Class & $K$ & \shortstack[l]{Materials\\per voxel} & \shortstack[l]{Materials\\per image} & Closure & Representative \\
\midrule
Two-basis & 2 & 2 & 2 & none needed & Alvarez and Macovski~\cite{alvarez1976} \\
Two-basis, more measurements & $>2$ & 2 & 2 & statistical weighting & Black et al.~\cite{blackprojection} \\
K-edge & $\geq3$ & $\geq3$ & $\geq3$ & none needed & Roessl and Proksa~\cite{roessl2007} \\
Three-material & 2 & 3 & 3 & conservation of volume or mass & Liu et al.~\cite{liu2009} \\
Multi-material & 2 & $\leq3$ & $>3$ & a list of triplets, one per voxel & Mendon\c{c}a et al.~\cite{mendonca2014} \\
Regularized multi-material & 2 & $\leq3$ & $>3$ & a penalty beside the list & Long and Fessler~\cite{long2014} \\
Single-energy & 1 & 2 & $>2$ & conservation of volume, one pair per voxel & Xue et al.~\cite{xue2021} \\
\bottomrule
\end{tabular}
\end{table*}

The third is how the answer is reported. Most methods return one estimate for each voxel. Bayesian methods
instead weigh several explanations of a measurement against prior
knowledge~\cite{lalonde2017bayes,soltani2024bayesvnc,jiang2025}.

% =============================================================================
\section{How the measurements are made}
\label{sec:acquisition}
% =============================================================================

\outline{Single energy, rapid kVp switching, dual source, dual layer, split filter, and photon counting, each as
a set of spectral responses $S_k(E)$. What sets the conditioning of the inversion: spectral separation, photon
budget, and consistent ray paths. The noise bound of the line integrals for each simulated system.}

\section{Projection-domain decomposition}
\label{sec:projection}

\outline{Calibration by polynomials and tables. Maximum likelihood for each ray. Two measurements and more than
two. Noise against the Cram\'er--Rao bound, bias at low counts, and the requirement of consistent rays.
Representative method written out in full.}

\section{Image-domain decomposition}
\label{sec:image}

\outline{Linear inversion in each voxel. Beam hardening and its precorrection. Noise amplification and the
correlation of the basis images. Regularized image-domain decomposition. Representative method.}

\section{One-step decomposition}
\label{sec:onestep}

\outline{Basis images estimated directly from the raw data. Penalized likelihood, surrogate and primal-dual
algorithms, cost and tuning. Representative method.}

\section{Beyond two materials}
\label{sec:multimaterial}

\outline{Three materials by the conservation of volume or mass. K-edge imaging with three or more measurements.
Multi-material decomposition with a list of triplets, the choice of the triplet, and its sensitivity.
Statistical and sparsity-regularized multi-material decomposition. Representative methods.}

\section{Single-energy decomposition}
\label{sec:single}

\outline{One equation for each voxel. Pairs of materials with the conservation of volume and a rule that selects
the pair. Optimization methods that group voxels. What one energy can and cannot determine.}

\section{Probabilistic decomposition}
\label{sec:probabilistic}

\outline{Several compositions can fit one measurement. Bayesian methods, priors on tissues, and posterior
sampling. What an uncertainty covers and how it is tested.}

\section{Learned decomposition}
\label{sec:learned}

\outline{Grouped by where the physics enters: nowhere (supervised networks), the loss (physics-informed
networks), the architecture (unrolled networks), or the output (methods fitted at test time). Grouped also by
the source of the labels: a dual-energy decomposition, simulation, or none. Accuracy inside and outside the
training data.}

\section{What decomposition is used for}
\label{sec:uses}

\outline{Virtual monoenergetic images, virtual non-contrast images, iodine maps, electron density and effective
atomic number, each as a linear combination of material maps, with its noise.}

\section{A common benchmark}
\label{sec:benchmark}

\outline{Every representative method on the same simulated scans of the same phantoms, with the same metrics.
The main table of the paper. Then the analyses that explain it.}

\section{Discussion}
\label{sec:discussion}

\outline{Paragraphs in the order of the sections. What the framework predicts about each class. Open problems.
Limitations of a simulated benchmark.}

\section{Conclusion}
\label{sec:conclusion}

\outline{The model, the three questions, and the benchmark, in three or four sentences.}

\section*{Data and code availability}
The notebook that produces every figure and number of this paper, with every method written out, is available
at \url{https://github.com/MolloiLab/MaterialDecompositionReview} and can be read as an interactive page at
\url{https://molloilab.github.io/MaterialDecompositionReview}. The simulator, \texttt{BasisSimulator.jl}, is
open source.

% =============================================================================
% References. Built by make_references.py from the verified entries of LITERATURE.md, in order of first
% citation. Never add an entry here by hand, and never from memory.
% =============================================================================
% BEGIN REFERENCES (generated by make_references.py from LITERATURE.md, do not edit)
\begin{thebibliography}{99}

\bibitem{alvarez1976}
R.~E. Alvarez and A. Macovski, ``Energy-selective reconstructions in X-ray computerised tomography,'' \emph{Phys. Med. Biol.}, vol.~21, no.~5, pp. 733--744, 1976, doi: 10.1088/0031-9155/21/5/002.

\bibitem{lehmann1981}
L.~A. Lehmann, R.~E. Alvarez, A. Macovski, W.~R. Brody, N.~J. Pelc, S.~J. Riederer, \emph{et al.}, ``Generalized image combinations in dual KVP digital radiography,'' \emph{Med. Phys.}, vol.~8, no.~5, pp. 659--667, 1981, doi: 10.1118/1.595025.

\bibitem{blackprojection}
D. Black, S. Naidu, H. Khodajou-Chokami, and S. Molloi, ``Unified projection-domain material decomposition for energy-integrating and photon-counting CT,'' 2026, arXiv:2610.30440.

\bibitem{roessl2007}
E. Roessl and R. Proksa, ``K-edge imaging in x-ray computed tomography using multi-bin photon counting detectors,'' \emph{Phys. Med. Biol.}, vol.~52, no.~15, pp. 4679--4696, 2007, doi: 10.1088/0031-9155/52/15/020.

\bibitem{mccollough2015}
C.~H. McCollough, S. Leng, L. Yu, and J.~G. Fletcher, ``Dual- and Multi-Energy CT: Principles, Technical Approaches, and Clinical Applications,'' \emph{Radiology}, vol.~276, no.~3, pp. 637--653, 2015, doi: 10.1148/radiol.2015142631.

\bibitem{rajendran2022}
K. Rajendran, M. Petersilka, A. Henning, E.~R. Shanblatt, B. Schmidt, T.~G. Flohr, \emph{et al.}, ``First Clinical Photon-counting Detector CT System: Technical Evaluation,'' \emph{Radiology}, vol.~303, no.~1, pp. 130--138, 2022, doi: 10.1148/radiol.212579.

\bibitem{mccollough2023}
C.~H. McCollough, K. Rajendran, S. Leng, L. Yu, J.~G. Fletcher, K. Stierstorfer, \emph{et al.}, ``The technical development of photon-counting detector CT,'' \emph{Eur. Radiol.}, vol.~33, no.~8, pp. 5321--5330, 2023, doi: 10.1007/s00330-023-09545-9.

\bibitem{maass2009}
C. Maa{\ss}, M. Baer, and M. Kachelrie{\ss}, ``Image-based dual energy CT using optimized precorrection functions: A practical new approach of material decomposition in image domain,'' \emph{Med. Phys.}, vol.~36, no.~8, pp. 3818--3829, 2009, doi: 10.1118/1.3157235.

\bibitem{long2014}
Y. Long and J.~A. Fessler, ``Multi-Material Decomposition Using Statistical Image Reconstruction for Spectral CT,'' \emph{IEEE Trans. Med. Imag.}, vol.~33, no.~8, pp. 1614--1626, 2014, doi: 10.1109/tmi.2014.2320284.

\bibitem{liu2009}
X. Liu, L. Yu, A.~N. Primak, and C.~H. McCollough, ``Quantitative imaging of element composition and mass fraction using dual-energy CT: Three-material decomposition,'' \emph{Med. Phys.}, vol.~36, no.~5, pp. 1602--1609, 2009, doi: 10.1118/1.3097632.

\bibitem{xue2021}
Y. Xue, W. Qin, C. Luo, P. Yang, Y. Jiang, T. Tsui, \emph{et al.}, ``Multi-Material Decomposition for Single Energy CT Using Material Sparsity Constraint,'' \emph{IEEE Trans. Med. Imag.}, vol.~40, no.~5, pp. 1303--1318, 2021, doi: 10.1109/tmi.2021.3051416.

\bibitem{mendonca2014}
P.~R.~S. Mendonca, P. Lamb, and D.~V. Sahani, ``A Flexible Method for Multi-Material Decomposition of Dual-Energy CT Images,'' \emph{IEEE Trans. Med. Imag.}, vol.~33, no.~1, pp. 99--116, 2014, doi: 10.1109/tmi.2013.2281719.

\bibitem{xue2017}
Y. Xue, R. Ruan, X. Hu, Y. Kuang, J. Wang, Y. Long, \emph{et al.}, ``Statistical image-domain multimaterial decomposition for dual-energy CT,'' \emph{Med. Phys.}, vol.~44, no.~3, pp. 886--901, 2017, doi: 10.1002/mp.12096.

\bibitem{bousse2024review}
A. Bousse, V.~S.~S. Kandarpa, S. Rit, A. Perelli, M. Li, G. Wang, \emph{et al.}, ``Systematic Review on Learning-Based Spectral CT,'' \emph{IEEE Trans. Radiat. Plasma Med. Sci.}, vol.~8, no.~2, pp. 113--137, 2024, doi: 10.1109/trpms.2023.3314131.

\bibitem{lalonde2017bayes}
A. Lalonde, E. B\"ar, and H. Bouchard, ``A Bayesian approach to solve proton stopping powers from noisy multi-energy CT data,'' \emph{Med. Phys.}, vol.~44, no.~10, pp. 5293--5302, 2017, doi: 10.1002/mp.12489.

\bibitem{soltani2024bayesvnc}
M. Beikali Soltani and H. Bouchard, ``Dual virtual non-contrast imaging: a Bayesian quantitative approach to determine radiotherapy quantities from contrast-enhanced DECT images,'' \emph{Phys. Med. Biol.}, vol.~69, no.~24, p. 245008, 2024, doi: 10.1088/1361-6560/ad965f.

\bibitem{jiang2025}
X. Jiang, G.~J. Gang, and J.~W. Stayman, ``Multi-Material Decomposition Using Spectral Diffusion Posterior Sampling,'' \emph{IEEE Trans. Biomed. Eng.}, vol.~72, no.~8, pp. 2447--2461, 2025, doi: 10.1109/tbme.2025.3543747.

\end{thebibliography}
% END REFERENCES

\end{document}
